\documentclass[10pt,a4paper]{amsart}
\usepackage[margin=19mm]{geometry}
\usepackage{amsmath,amssymb,mathtools}
\usepackage[scr=rsfs,cal=euler]{mathalfa}
\usepackage{xfrac}
\usepackage[unicode,hidelinks,bookmarks=false]{hyperref}
\newcommand{\ie}{i.e.\ }
\newcommand{\cf}{cf.\ }
\newcommand{\resp}{respectively\ }
\newcommand{\Subsection}{\subsection}
\providecommand{\nobreakdash}{\nobreak\hbox{-}}
\setlength{\emergencystretch}{2em}
\multlinegap0pt
\usepackage{bm}
\usepackage{bbm}
\usepackage{dsfont}
\usepackage{xspace}
\usepackage{cancel}
\usepackage{mathtools}
\usepackage{amsthm}
\makeatletter
\@ifundefined{theorem}{\newtheorem{theorem}{Theorem}}{}
\@ifundefined{lemma}{\newtheorem{lemma}[theorem]{Lemma}}{}
\@ifundefined{proposition}{\newtheorem{proposition}[theorem]{Proposition}}{}
\makeatother
\newcommand{\R}{\mathbb{R}}
\title[Connected components and topological ends of stationary plan]{Connected components and topological ends of stationary planar forests}
\author{Tom Garcia-Sanchez}
\date{Source: arXiv:2604.04672v2 (CC BY 4.0)}
\begin{document}
\maketitle

\noindent\emph{Source and licence.} Connected components and topological ends of stationary planar forests. Version arXiv:2604.04672v2 (CC BY 4.0), distributed under the Creative Commons Attribution 4.0 International licence, as recorded in the arXiv licence metadata for this exact version and shown on the abstract page https://arxiv.org/abs/2604.04672v2. Full rights notice: licenses/arxiv-2604.04672-CC-BY-4.0.txt.\par\smallskip

\noindent\textit{Editorial scope.} Counted unit: the Proposition whose source label is \texttt{prop\_arc\_diff} in the article (source0.tex lines 875--882, source label \texttt{prop\_arc\_diff}), delivered together with the article's own complete original proof of it (source0.tex lines 884--923, one \texttt{proof} environment of 40 lines). No mathematics is written, repaired or re-derived: statement and proof are the article's own text line for line. Required context retained uncounted: lemma\_adapted\_homeo (lines 832--838). Every definition, display and label of those lines is retained unchanged. Omitted from this package and not redistributed: 1--831, 839--874, 883--883, 924--1096. Nothing inside a retained excerpt is omitted or altered: every retained body is delivered exactly as the article's lines are written, so no omission receipt of any kind accompanies this package. Citations: the retained excerpts carry no citation command at all, so no bibliography entry of the article is transcribed and no reference is redistributed. Reference adaptation: none was needed and none was made; every reference inside the delivered unit resolves inside it, so no unresolved reference or citation is produced. Printed slips kept literal: no printed word, symbol, hypothesis or number of the counted statement or of its proof was altered. Layout and class: the article's own class is \texttt{article} with packages including geometry, fontenc, lmodern, amssymb, mathtools, stmaryrd, graphicx, pgfplots, mathrsfs, dsfont, amsthm, enumerate, color, verbatim; the wrapper is the lane's pinned \texttt{amsart} editorial wrapper at 10pt on A4, which restates the theorem-like environments the retained text needs and adds the article's own macro definitions that the retained text needs (\texttt{\textbackslash R}), with extra wrapper packages (bm, bbm, dsfont, xspace, cancel, mathtools). The delivered numbering is the unit's own. Assets: no retained span contains a figure environment, a graphics-inclusion command, a PDF-page inclusion, a table or a TikZ picture; no image file is copied into this package, the package declares no asset dependency and no image file is redistributed with it. Licence: version arXiv:2604.04672v2 of this article is distributed under the Creative Commons Attribution 4.0 International licence, as recorded in the arXiv licence metadata for this exact version and shown on the abstract page https://arxiv.org/abs/2604.04672v2. A full rights notice is delivered at licenses/arxiv-2604.04672-CC-BY-4.0.txt. Characters: every retained excerpt is pure ASCII, so the delivered engine, XeLaTeX with its default Latin Modern text font, reports no missing character.
\begin{lemma}
    \label{lemma_adapted_homeo}
    Fix $\gamma\in\Gamma$ and an homeomorphism $\Phi:\R\to\gamma$. For each $\sigma\in\Gamma\setminus\{\gamma\}$, there exists a homeomorphism $\phi_{\sigma}$ and a non-empty bounded open interval $I_{\sigma}\coloneqq (a_\sigma,b_\sigma)\subset\R$ such that
        \[\text{$\Phi=\phi_\sigma$ on $\R\setminus I_\sigma$},\quad \phi_\sigma(I_\sigma)=\sigma\setminus\gamma,\quad\text{and}\quad \Phi(I_\sigma)=\gamma\setminus\sigma.\]
    Then, one has additionally that for all distinct $\sigma,\sigma'\in\Gamma\setminus\{\gamma\}$,
        \[(\partial I_{\sigma'})\setminus\overline{I_{\sigma}}\subset \overline{I_{\sigma'|\sigma}}\quad\text{where}\quad I_{\sigma'|\sigma}\coloneqq\phi_{\sigma'}^{-1}(\sigma'\setminus\sigma).\]
\end{lemma}

\begin{proposition}
    \label{prop_arc_diff}
    For all distinct $\gamma,\gamma'\in\Gamma$,
        \[J(\gamma,\gamma')\coloneqq\overline{\gamma\triangle\gamma'}\]
    is a Jordan curve. Furthermore, for any $\gamma\in\Gamma$ and non-empty finite subset $\Gamma_0\subset\Gamma\setminus\{\gamma\}$,
        \[\gamma\setminus\bigcup_{\sigma\in\Gamma_0}\sigma\]
    is a non-empty portion of $\gamma$, and for all distinct $\gamma,\sigma,\sigma'\in\Gamma$, the difference $J(\gamma,\sigma')\setminus J(\gamma,\sigma)$ is a connected subset of $J(\sigma,\sigma')$.
\end{proposition}

\begin{proof}
    Fix $\gamma\in\Gamma$ and an homeomorphism $\Phi:\R\to\gamma$. Let 
        \[(\phi_\sigma)_{\sigma\in\Gamma\setminus\{\gamma\}}\quad\text{and}\quad(I_\sigma)_{\sigma\in\Gamma\setminus\{\gamma\}}=[(a_\sigma,b_\sigma)]_{\sigma\in\Gamma\setminus\{\gamma\}}\]
    be the homeomorphisms and non-empty bounded open intervals given by Lemma~\ref{lemma_adapted_homeo}. Let $\gamma'\in\Gamma\setminus\{\gamma\}$. Since by construction $\Phi(I_{\gamma'})=\gamma\setminus\gamma'$ and $\phi_{\gamma'}(I_{\gamma'})=\gamma'\setminus\gamma$ are disjoint, and $\Phi=\phi_{\gamma'}$ on the two distinct points of $\partial I_{\gamma'}$, then
        \[J(\gamma,\gamma')=\Phi(\overline{I_{\gamma'}})\cup\phi_{\gamma'}(\overline{I_{\gamma'}})\]
    is a Jordan curve.

    Now, let us fix a non-empty finite subset $\Gamma_0\subset\Gamma\setminus\{\gamma\}$ and show that $\gamma\setminus\bigcup_{\sigma\in\Gamma_0}\sigma$ is a non-empty portion of $\gamma$. By construction,
        \[\gamma\setminus\bigcup_{\sigma\in\Gamma_0}\sigma=\Phi(I_{\Gamma_0})\quad\text{where}\quad I_{\Gamma_0}\coloneqq\bigcap_{\sigma\in\Gamma_0}I_\sigma.\]
    Let $\sigma,\sigma'\in\Gamma_0$ be such that $b_{\sigma'}=\max_{\nu\in\Gamma_0} b_\nu$ and $a_{\sigma}=\min_{\nu\in\Gamma_0} a_\nu$. Then, $a_{\sigma}\leq b_{\sigma'}$ and
        \[I_{\Gamma_0}=(a_{\sigma}, b_{\sigma'}).\]
    It suffices therefore to show that $a_{\sigma}<b_{\sigma'}$. Assume by contradiction that $a_{\sigma}=b_{\sigma'}$. Then $I_{\sigma}$ and $I_{\sigma'}$ must be disjoint. Fix $t\in I_{\sigma}$. By definition, $\Phi(t)\notin\sigma$, and $\phi_{\sigma'}(t)=\Phi(t)$ since $t\notin I_{\sigma'}$. This gives $\phi_{\sigma'}(t)\in\sigma'\setminus\sigma$ and hence
        \[t\in I_{\sigma'|\sigma}.\]
    Additionally, both $a_{\sigma}$ and $b_{\sigma}$ belong to $\R\setminus(I_{\sigma}\cup I_{\sigma'})$, so that $\phi_{\sigma'}(a_{\sigma})=\Phi(a_{\sigma})=\phi_{\sigma}(a_{\sigma})\in\sigma$ and $\phi_{\sigma'}(b_{\sigma})=\Phi(b_{\sigma})=\phi_{\sigma}(b_{\sigma})\in\sigma$, which implies that
        \[a_{\sigma}\notin I_{\sigma'|\sigma}\quad\text{and}\quad b_{\sigma}\notin I_{\sigma'|\sigma}.\]
    Together, it shows that $I_{\sigma'|\sigma}$ must contain $t\in(a_\sigma,b_\sigma)$ but not $a_\sigma$ nor $b_\sigma$. Since $\sigma'\setminus\sigma$ is a portion of $\sigma'$, $I_{\sigma'|\sigma}=\phi_{\sigma'}^{-1}(\sigma'\setminus\sigma)$ must be an open interval, thus one must have
        \[I_{\sigma'|\sigma}\subset I_\sigma.\]
    But then, applying Lemma~\ref{lemma_adapted_homeo}, it comes $(\partial I_{\sigma'})\setminus\overline{I_\sigma}\subset\overline{I_{\sigma'|\sigma}}\subset\overline{I_\sigma}$, i.e., $\partial I_{\sigma'}\subset\overline{I_\sigma}$, which is absurd since $b_{\sigma'}\in\partial I_{\sigma'}$ but $b_{\sigma'}\notin\partial I_\sigma$ as $b_{\sigma'}>a_{\sigma'}\geq b_\sigma$.

    It now remains to prove the last assertion. Fix distinct $\sigma,\sigma'\in\Gamma\setminus\{\gamma\}$. Let us prove that $J(\gamma,\sigma')\setminus J(\gamma,\sigma)$ is a connected subset of $J(\sigma,\sigma')$. First, observe that by construction, one can write
        \[J(\gamma,\sigma')=\Phi(\overline{I_{\sigma'}})\cup \phi_{\sigma'}(I_{\sigma'})\quad\text{and}\quad J(\gamma,\sigma)=\Phi(\overline{I_{\sigma}})\cup \phi_{\sigma}(I_{\sigma}).\]
    Then, taking the difference, it comes
        \begin{equation}
            \label{eq_Jordan_diff}
            \begin{split}
                J(\gamma,\sigma')\setminus J(\gamma,\sigma)&=\Phi(\overline{I_{\sigma'}})\setminus[\Phi(\overline{I_{\sigma}})\cup \phi_{\sigma}(I_{\sigma})]\cup\phi_{\sigma'}(I_{\sigma'})\setminus[\Phi(\overline{I_\sigma})\cup\phi_{\sigma}(I_{\sigma})]\\
                &=\Phi(\overline{{I_{\sigma'}}}\setminus\overline{I_\sigma})\cup\phi_{\sigma'}(I_{\sigma'})\setminus\phi_{\sigma}(I_{\sigma})\\
                &=\Phi(\overline{{I_{\sigma'}}}\setminus\overline{I_\sigma})\cup\phi_{\sigma'}(I_{\sigma'}\cap I_{\sigma'|\sigma})
            \end{split}
        \end{equation}
    where using the definition of $I_\sigma$ and $I_{\sigma'}$, the first equality follows from $\phi_\sigma(I_\sigma)\cap\Phi(\overline{I_{\sigma'}})\subset\phi_\sigma(I_\sigma)\cap\gamma=\emptyset$ as well as $\phi_{\sigma'}(I_{\sigma'})\cap\Phi(\overline{I_\sigma})\subset\phi_{\sigma'}(I_{\sigma'})\cap\gamma=\emptyset$, and $\phi_\sigma(\R\setminus I_{\sigma})\cap\phi_{\sigma'}(I_{\sigma'})\subset\gamma\cap\phi_{\sigma'}(I_{\sigma'})=\emptyset$ yields the last line. Now, by construction, $\phi_{\sigma'}(I_{\sigma'}\cap I_{\sigma'|\sigma})\subset\phi_{\sigma'}(I_{\sigma'|\sigma})=\sigma'\setminus\sigma$, and
        \[\Phi(\overline{{I_{\sigma'}}}\setminus\overline{I_\sigma})\subset\overline{\Phi(I_{\sigma'})\setminus\Phi(I_\sigma)}\subset\overline{(\gamma\setminus\sigma')\setminus(\gamma\setminus\sigma)}\subset\overline{\sigma\setminus\sigma'}.\]
    Injected in~\eqref{eq_Jordan_diff}, it shows
        \[J(\gamma,\sigma')\setminus J(\gamma,\sigma)\subset\overline{\sigma\setminus\sigma'}\cup\sigma\setminus\sigma'\subset \overline{\sigma\triangle\sigma'}=J(\sigma,\sigma').\]
    It remains to check connectedness. To that end, observe that since $I_{\sigma'}\cap I_{\sigma'|\sigma}$ is an open interval,
        \[\phi_{\sigma'}(I_{\sigma'}\cap I_{\sigma'|\sigma})\]
    is connected. Then, observe that although $\overline{{I_{\sigma'}}}\setminus\overline{I_\sigma}$ may be disconnected, each of its non-empty components must contain a point of $(\partial I_{\sigma'})\setminus\overline{I_\sigma}$. Finally, since
        \[(\partial I_{\sigma'})\setminus\overline{I_\sigma}\subset\overline{I_{\sigma'|\sigma}}\cap \overline{I_{\sigma'}}=\overline{I_{\sigma'}\cap I_{\sigma'|\sigma}}\]
    from Lemma~\ref{lemma_adapted_homeo}, each component of $\Phi(\overline{{I_{\sigma'}}}\setminus\overline{I_\sigma})$ connects to $\phi_{\sigma'}(I_{\sigma'}\cap I_{\sigma'|\sigma})$, so that injecting in~\eqref{eq_Jordan_diff}, it comes $J(\gamma,\sigma')\setminus J(\gamma,\sigma)$ is connected.
\end{proof}

\end{document}
